\section{Day 3: Random Variables (Sep.\ 16, 2026)}
Let $(\Omega, \SA, \PP)$ be a probability space, and let $(S, \mathcal S)$ be a measurable space. A random variable is a \textit{measurable function} $X \colon \Omega \to S$ (i.e., for all $U \in \mathcal S$, $X\inv(U) \in \SA$). For example, consider $(\Omega, \SA, \PP) = ([0, 1], \SB([0, 1]), dx)$ (i.e., the Borel $\sigma$-algebra with the Lebesgue measure). Let $S = \{H, T\}$ be the outcomes of a coinflip; we may define a random variable $X$ as follows,
\[ X(\omega) = \begin{cases} H, &\omega \geq \frac{1}{2}, \\ T, &\omega < \frac{1}{2}. \end{cases} \]
Note that this is not the only way to define a fair coin; we could consider a random variable $\tilde X$ that sends $\omega$ to $H$ if $\min\{\abs{\omega}, \abs{1 - \omega}\} < \frac{1}{4}$, or any other fancy examples that one could come up with.

For future reference, we will often omit prescribing the probability space $(\Omega, \SA, \PP)$ if it is obvious.

\begin{definition} \label{defn:distribution}
    The law/distribution of a random variable $X$ is the probability measure on $(S, \mathcal S)$ given by $\PP_X(U) = \PP(X \in U) = \PP(X\inv(U))$.
\end{definition}

You may recognize this as the pushforward measure; though probabilists will often not refer to it as such. Following our example from earlier,
\[ \PP(X = H) = \PP(X\inv(H)) = \PP\left(\left[\frac{1}{2}, 1\right]\right) = \frac{1}{2}. \]

\begin{definition} \label{defn:equal_in_distribution}
    Two random variables $X, Y$ are equal in law/distribution if $\PP_X = \PP_Y$. We write $\smash{X \stackrel{d}{=} Y}$.
\end{definition}

For example, $\smash{X \stackrel{d}{=} \tilde X}$, as both are fair coins.

\begin{example}[Durrett, 1.2.2]
    Suppose we have a CDF $F \colon \RR \to [0, 1]$. How can we define a random variable $X$ with CDF $F$? For one, we may set $(\Omega, \SA, \PP) = (\RR, \SB(\RR), \PP_\SA)$, with $X(\omega) = \omega$. Alternatively, we may set $(\Omega, \SA, \PP) = ([0, 1], \SB([0, 1]), dx)$, with
    \[ X(\omega) = \inf \{s \in \RR : F(s) \geq \omega\}. \]
    Essentially, we want $\PP(X \leq t) = F(t)$. Observe that
    \begin{align*}
        X \leq t &\iff \inf\{s \in \RR : F(s) \geq \omega\} \leq t \\
        &\iff \lim_{s \to t^+} F(s) \geq \omega.
    \end{align*}
    Note that we're taking the limit from above here, since $F$ is not necessarily continuous; though, it is right-continuous. In this manner, the above is equivalent to $F(t) \geq \omega$, so
    \[ \PP(X \leq t) = dx(\omega \in [0, F(t)]) = F(t). \]
\end{example}

\begin{definition} \label{defn:sigma_algebra_generated_by_rv}
    Suppose $X \colon (\Omega, \SA, \PP) \to (S, \mathcal S)$ is a random variable. We write $\sigma(X) = \{X\inv(A) \colon A \in \mathcal S\}$ for the $\sigma$-algebra generated by $X$.
\end{definition}
In particular, $\sigma(X)$ is the smallest $\sigma$-algebra such that $X$ is measurable. As an example, consider $X$ to be the coinflip from earlier, with heads for $\omega$ above one-half and tails otherwise; then
\[ \sigma(X) = \left\{\emptyset, [0, 1], \left[\frac{1}{2}, 1\right], \left[0, \frac{1}{2}\right)\right\}. \]

\newpage
\begin{lemma}[Panchenko 1.4] \label{lem:transport_of_measure}
    Suppose $X, Y$ are random variables on $(\Omega, \SA, \PP)$, taking values in $(S, \mathcal S)$ and $(\RR, \SB(\RR))$ respectively. Then $Y$ is $\sigma(X)$-measurable, written ``$Y \in \sigma(X)$'', if and only if $Y = g(X)$ for some measurable function $g \colon (S, \mathcal S) \to (\RR, \SB(\RR))$.
\end{lemma}
\begin{proof}
    \begin{itemize}
        \item[($\Rightarrow$)] For a natural $n \in \NN$ and integer $k \in \ZZ$, define
        \[ A_{n,k} = Y\inv \left(\left[\frac{k}{2^n}, \frac{k+1}{2^n}\right]\right) \in \sigma(X). \]
        Then $(A_{n,k} : k \in \ZZ)$ is a $\sigma(X)$-measurable partition of $\Omega$, so $A_{n,k} = X\inv(B_{n,k})$ for some disjoint set $(B_{n,k} : k \in \ZZ)$ in $\mathcal S$. Next, set $g_n \colon S \to \RR$ given by
        \[ g_n(x) = \sum_{k \in \ZZ} \frac{k}{2^n} \mathbf 1(x \in B_{n,k}). \]
        Clearly, the $g_n$'s are increasing, i.e., $g_n \leq g_{n+1}$ for all $n$. Denote $g = \lim_{n \to \infty} g_n$ (note that the limit exists by monotonicity, and the limit of measurable functions is measurable). Finally, we get that $g(X) = Y$.
        \item[($\Leftarrow$)] For any borel set $B \subset \RR$, let $B' = g\inv(B) \in \SB$; then
        \[ \{Y = g(X) \in B\} = \{X \in g\inv(B) = B'\} = X\inv(B') \in \sigma(X). \qedhere \]
    \end{itemize}
\end{proof}

We now talk about independence. Let $(\Omega, \SA, \PP)$ be a probability space.
\begin{definition} \label{defn:independence}
    The events $A_1, \dots, A_n \in \SA$ are \textit{independent} if, for all $I \subset \{1, \dots, n\}$,
    \[ \PP\left(\bigcap_{i \in I} A_i\right) = \prod_{i \in I} \PP(A_i). \]
\end{definition} \vspace{-8pt}
As an example, consider a tetrahedral die colored white, blue, red, and rainbow (i.e., all three colors) on each of its sides. The colors that are rolled are examples of pairwise independent random variables, but are not independent altogether (cf.\ Example 1.2.2 in Panchenko).

\begin{definition} \label{defn:independence_for_collections}
    Collections of events $\SA_1, \dots, \SA_n$ are independent if, for all $A_1 \in \SA_1$, $A_2 \in \SA_2$, etc., we have that $A_1, \dots, A_n$ are independent.
\end{definition}

\begin{definition} \label{defn:independence_for_rvs}
    Random variables $X_1, \dots, X_n$ are independent if $\sigma(X_1), \dots, \sigma(X_n)$ are independent.
\end{definition}
The above is equivalent to asking that, for all measurable sets $A_1, \dots, A_n$ in the codomains of $X_1, \dots, X_n$ respectively, we have that
\[ \PP(X_i \in A_i : \forall i = 1, \dots, n) = \prod_{i=1}^n \PP(X_i \in A_i). \]
\vspace{-16pt}
\begin{lemma}[Panchenko 1.6] \label{lem:independence_for_pi_systems}
    If $\SCC_1, \dots, \SCC_n$ are independent $\pi$-systems, then $\sigma(\SCC_1), \dots, \sigma(\SCC_n)$ are independent.
\end{lemma}
\begin{proof}
    By induction, we want to show that $\sigma(\SCC_1), \dots, \sigma(\SCC_n)$ are independent. We want to apply the Dynkin $\pi$-$\lambda$ theorem. Let $\mathcal L = \{ A : \{A\}, \SCC_2, \dots, \SCC_n \text{ are independent}\}$. We have four things to check, namely, that $\SCC_1 \subset \mathcal L$ by assumption, $\mathcal L$ is non-empty, admits complements, and is closed under countable unions. Respectively, this is true by assumption, that $\Omega \in \mathcal L$, that, for $A \in \mathcal L$, we have, for $A_2, \dots, A_n$ in $\SCC_2, \dots, \SCC_n$ respectively,
    \[ \PP(A^c \cap A_2 \cap \dots \cap A_n) = \PP(A_2 \cap \dots \cap A_n) - \PP(A \cap A_2 \cap \dots \cap A_n). \]
    The last part is left as an exercise.
\end{proof}

\newpage
\begin{definition} \label{defn:joint_cdf}
    For real-valued random variables $X_1, \dots, X_n$ taking values in $(\RR, \SB(\RR))$, the \textit{joint CDF} is the function
    \[ F_{(X_1, \dots, X_n)} = F \colon \RR^n \to [0, 1] \]
    given by $F(x_1, \dots, x_n) = \PP(X_i \leq x_i : \forall i = 1, \dots, n)$.
\end{definition}
If $X_1, \dots, X_n$ are independent, then $F_{(X_1, \dots, X_n)} = F_{X_1} \dots F_{X_n}$.
\begin{proof}
    The forwards direction is easy. The converse direction follows from the proposition to check independence of $X_1, \dots, X_n$. It is enough to check independence of the $\pi$-systems
    \[ \left(X_i\inv (-\infty, x] : x \in \RR\right) \]
    for $i = 1, \dots, n$, since they generate $\sigma(X_i)$.
\end{proof}

\begin{lemma}[Grouping Lemma] \label{lem:grouping_lemma}
    If $\{\SA_{i,j} : i = 1, \dots, n; j = 1, \dots, \alpha(i)\}$ is a collection of independent collections of sets, then
    \[ \left\{ \sigma\left(\SA_{i,1} \cup \dots \cup \SA_{i,\alpha(i)} : i = 1, \dots, n \right)\right\} \]
    are independent.
\end{lemma}
\begin{proof}
    Left as an exercise.
\end{proof}

For example, if $X_1, \dots, X_4$ are independent random variables, then $\cos(X_1 X_2)$ and $\tanh(X_3 X_4)$ are independent.\footnote{however the fuck we got here}

{\color{airforceblue}``If you got a 7/100 on the midterm, and your friend got 9/100 on the midterm... what's the first thing you're going to ask me? What's the average?''}

We now go over expected value.
\begin{definition} \label{defn:expected_value}
    If $X \colon \Omega \to \RR$ is a real-valued random variable, then
    \[ \EE X = \int_{\Omega} X(\omega) \, d\PP(\omega). \]
    For the above definition to exist, we need $\EE \abs{X} = \int_\Omega \abs{X}\!(\omega) \, d\PP(\omega) < \infty$.
\end{definition}

\begin{fact} \label{fact:expected_value}
    If $\PP_X$ is the law of $X$, then
    \[ \EE X = \int_\RR x \, d \PP_X(x). \]
\end{fact}

This is immediately obtained from \ref{defn:expected_value} by change-of-variable. In general,
\[ \EE g(X) = \int_\RR g(x) \, d\PP_X(x) = \int_\RR x \, d\PP_{g(X)}(x). \]